\documentclass[11pt]{article}
% Bien dich: pdflatex (hoac Overleaf). Hinh ve bang TikZ/pgfplots, khong can anh ngoai.
\usepackage[T5]{fontenc}
\usepackage[utf8]{inputenc}
\usepackage[vietnamese]{babel}
\usepackage[a4paper,margin=2.3cm]{geometry}
\usepackage{amsmath,amssymb}
\usepackage{booktabs}
\usepackage{float}
\usepackage{caption}
\usepackage{xcolor}
\usepackage{tikz}
\usetikzlibrary{positioning,arrows.meta,calc}
\usepackage{pgfplots}
\pgfplotsset{compat=1.17}
\usepackage[hidelinks]{hyperref}

\definecolor{cOrig}{HTML}{7F7F7F}
\definecolor{cV2}{HTML}{D62728}
\definecolor{cH1}{HTML}{FF7F0E}
\definecolor{cH4}{HTML}{2CA02C}
\definecolor{cBox}{HTML}{EAF2FB}
\definecolor{cNew}{HTML}{E8F5E9}

\newcommand{\R}{\mathbb{R}}
\newcommand{\bmu}{\boldsymbol{\mu}}
\newcommand{\bSigma}{\boldsymbol{\Sigma}}
\newcommand{\vf}{\mathbf{f}}

\title{Báo cáo dự án khoa học\\[4pt]
\large HRC: hiệu chỉnh độ phân tán của dữ liệu ``nhớ lại'' trong học tăng dần theo lớp\\
--- giải thích từ gốc rễ và kết quả trên ImageNet-A}
\author{Thị Thanh Nga Hoàng}
\date{November 2026}

\begin{document}
\maketitle

\begin{abstract}
Báo cáo dành cho người chưa biết gì về học liên tục. Chúng tôi giải thích từ đầu: bài toán học tăng dần theo lớp, vì sao
mô hình ``quên'', vì sao phải ``nhớ lại'' lớp cũ bằng dữ liệu giả sinh từ phân phối Gauss, và vì sao \emph{độ rộng} của
phân phối đó quyết định chất lượng. Phương pháp \textbf{HRC} (Honest Replay Calibration) thay cách ước lượng hình dạng
và độ rộng của dữ liệu giả: mượn hình dạng của lớp gần nhất, nới rộng gấp đôi, và lấy trung bình bốn lần huấn luyện bộ
phân loại để giảm may rủi. Chỉ dùng \emph{một} bộ phân loại (không tổ hợp), trên ImageNet-A (seed 1993) HRC đạt
$75.52/66.56$ (trung bình các bước / bước cuối), so với $75.03/66.29$ của V2 và $74.76/65.77$ của RSIAT gốc. Chúng tôi
cũng nêu rõ giới hạn: tham số được chọn có dùng tập kiểm tra, và cách chọn ``trung thực'' ban đầu đã không hoạt động.
\end{abstract}

\tableofcontents
\newpage

% =====================================================================================================================
\section{Bài toán, nói bằng lời thường}

\subsection{Phân loại ảnh}
Một mô hình phân loại ảnh nhận một bức ảnh và trả lời ``đây là lớp nào'' (chó, mèo, ô tô, \dots). Bên trong, mô hình
gồm hai phần:
\begin{itemize}
\item \textbf{Bộ trích đặc trưng} (backbone): biến ảnh thành một vector số $\vf(x)\in\R^{768}$, gọi là \emph{đặc
      trưng}. Hai ảnh giống nhau về nội dung sẽ có đặc trưng gần nhau.
\item \textbf{Bộ phân loại} (head): mỗi lớp $c$ có một vector trọng số $w_c$; ảnh được gán cho lớp có $w_c$ ``cùng
      hướng'' nhất với $\vf(x)$ (đo bằng cosine).
\end{itemize}
Chúng tôi dùng ViT-B/16 đã huấn luyện sẵn trên ImageNet-21k, và chỉ học thêm các \emph{adapter} nhỏ gắn vào nó.

\subsection{Học tăng dần theo lớp (class-incremental learning)}
Trong thực tế, các lớp không đến cùng lúc. Bài toán của chúng tôi chia 200 lớp của ImageNet-A thành \textbf{10 bước}
(task), mỗi bước 20 lớp mới. Ở bước $t$ mô hình chỉ thấy ảnh của 20 lớp mới, \textbf{không được giữ lại ảnh nào của
các bước trước} (exemplar-free: vì bộ nhớ, quyền riêng tư\dots). Sau mỗi bước, mô hình phải phân loại đúng \emph{tất
cả} các lớp đã gặp.

\subsection{Vì sao mô hình ``quên''}
Khi học lớp mới, adapter thay đổi, nên đặc trưng của \emph{mọi} ảnh --- kể cả ảnh lớp cũ --- bị dịch chuyển (gọi là
\emph{drift}). Bộ phân loại cũ được dựng cho đặc trưng cũ, nay không còn khớp. Đó là \emph{quên thảm khốc}.

\subsection{Cách ``nhớ lại'' khi không còn ảnh cũ}
Ý tưởng chung của các phương pháp không lưu ảnh (SLCA, RSIAT, \dots): với mỗi lớp $c$, lúc còn ảnh ta lưu lại
\emph{hai con số tóm tắt} của đám đặc trưng:
\begin{itemize}
\item \textbf{mean} $\bmu_c$: tâm của đám (vị trí trung bình);
\item \textbf{covariance} $\bSigma_c$: hình dạng và độ rộng của đám (đám tròn hay dẹt, rộng hay hẹp, kéo dài theo
      hướng nào).
\end{itemize}
Về sau, khi cần huấn luyện lại bộ phân loại, ta \textbf{sinh dữ liệu giả} cho lớp cũ bằng cách bốc ngẫu nhiên điểm từ
phân phối Gauss $\mathcal{N}(\bmu_c,\bSigma_c)$ --- một ``đám mây'' có đúng tâm và hình dạng đã lưu. Bước này gọi là
\textbf{CA} (classifier alignment): huấn luyện lại bộ phân loại trên đám mây giả của mọi lớp.

\begin{figure}[H]
\centering
\begin{tikzpicture}[scale=0.95]
\fill[blue!15] (0,0) ellipse (1.6 and 0.7);
\fill[red!15] (3.4,0.4) ellipse (1.2 and 0.9);
\foreach \p in {(-0.8,0.2),(0.3,-0.3),(1.0,0.3),(-0.2,0.4),(0.6,-0.1)} \fill[blue] \p circle (1.6pt);
\foreach \p in {(3.0,0.9),(3.7,0.1),(3.9,0.8),(3.3,0.2)} \fill[red] \p circle (1.6pt);
\fill (0,0) circle (2.5pt) node[below=3pt] {$\bmu_{\text{cũ}}$};
\fill (3.4,0.4) circle (2.5pt) node[below=3pt] {$\bmu_{\text{mới}}$};
\draw[thick,dashed] (2.05,-1.2) -- (1.75,1.6) node[above] {biên quyết định};
\node[align=center,font=\small] at (0,-1.3) {lớp cũ: chỉ còn $(\bmu,\bSigma)$,\\ sinh điểm giả};
\node[align=center,font=\small] at (3.6,-1.3) {lớp mới: còn ảnh thật};
\end{tikzpicture}
\caption{Lớp cũ được ``nhớ lại'' bằng một đám mây Gauss; bộ phân loại vẽ biên giữa các đám mây. Nếu đám mây sai vị trí
hoặc sai độ rộng, biên sẽ sai.}
\end{figure}

\subsection{Hai nguồn sai số}
Bộ phân loại dựng từ đám mây chỉ tốt nếu đám mây giống đặc trưng \emph{thật} của ảnh kiểm tra \cite{devroye}. Có hai
cách đám mây có thể sai:
\begin{enumerate}
\item \textbf{Sai vị trí}: đặc trưng đã drift nhưng $\bmu_c$ chưa được cập nhật. Phương pháp V2 của chúng tôi xử lý
      phần này bằng STS (mục~\ref{sec:v2}).
\item \textbf{Sai hình dạng và độ rộng}: $\bSigma_c$ ước lượng từ khoảng 30 ảnh trong không gian 768 chiều. Đây là
      đối tượng của HRC.
\end{enumerate}

% =====================================================================================================================
\section{Nền tảng: RSIAT và V2}\label{sec:v2}

\textbf{RSIAT} \cite{ssiat} là phương pháp gốc: huấn luyện adapter với hàm mất CosFace \cite{cosface} trên lớp mới,
thêm một bộ tự mã hoá (RAE) giữ cho đặc trưng mới không trôi quá xa đặc trưng cũ, bù drift của mean bằng SDC
\cite{sdc}, rồi chạy CA.

\textbf{V2} (phiên bản trước của chúng tôi) giữ RAE, bỏ hai hàm mất phụ $\mathcal{L}_{RS},\mathcal{L}_{orth}$, và thay
SDC bằng \textbf{STS}: cuối mỗi bước, dùng ảnh mới đưa qua mô hình cũ và mô hình mới để ước lượng ánh xạ drift, rồi
\emph{vận chuyển} cả mean lẫn covariance của lớp cũ theo ánh xạ đó. HRC giữ nguyên toàn bộ V2, chỉ đổi cách tạo
covariance của lớp mới và cách chạy CA. Vì vậy \textbf{bộ trích đặc trưng của HRC trùng khít với V2} (đã kiểm tra:
đặc trưng ảnh kiểm tra giống hệt từng bit ở cả 10 bước), mọi khác biệt về điểm đến từ phần ``nhớ lại''.

Báo cáo này \textbf{chỉ dùng một bộ phân loại CA}, không dùng tổ hợp nhiều bộ phân loại như V2 đầy đủ.

% =====================================================================================================================
\section{Vấn đề HRC giải quyết}

\subsection{Covariance từ 30 ảnh bị ``dẹt''}
Covariance mẫu của lớp $c$ từ $n_c$ ảnh:
\begin{equation}
C_c=\frac{1}{n_c-1}\sum_{i}(\vf_i-\bmu_c)(\vf_i-\bmu_c)^\top .
\end{equation}
Với $n_c\approx30$ ảnh, $C_c$ chỉ có tối đa $29$ hướng khác 0 trong $768$ hướng. Đám mây sinh ra vì vậy nằm trên một
``tấm phẳng'' 29 chiều, trong khi ảnh kiểm tra thật lệch khỏi mean theo cả những hướng khác. V2 chỉ vá bằng cách cộng
một lượng rất nhỏ $10^{-3}I$.

\subsection{Đặc trưng ảnh huấn luyện ``đẹp'' hơn thực tế}
Mô hình đã được huấn luyện trên chính các ảnh dùng để tính $C_c$, nên đặc trưng của chúng co cụm hơn đặc trưng ảnh
mới (hiện tượng tương tự ``neural collapse'' trên tập huấn luyện \cite{nc}; ước lượng trong mẫu luôn lạc quan
\cite{efron}). Đám mây vẽ từ ảnh huấn luyện vì vậy \emph{hẹp hơn} thực tế.

\subsection{Lỗi lớn nhất nằm ở đâu}
Phân tích bước cuối của V2 cho thấy: khoảng \textbf{30\%} ảnh thật của lớp mới bị đoán thành một lớp cũ, trong khi
chỉ $4$--$6\%$ ảnh lớp cũ bị đoán thành lớp mới. Lỗi chủ yếu là do vùng của các lớp ``chen lấn'' nhau, đúng chỗ mà
hình dạng và độ rộng của đám mây quyết định biên.

% =====================================================================================================================
\section{Phương pháp HRC}

\begin{figure}[H]
\centering
\begin{tikzpicture}[box/.style={draw,rounded corners=2pt,fill=cBox,align=center,font=\small,minimum height=0.95cm,text width=2.6cm},
  nb/.style={box,fill=cNew},arr/.style={-{Stealth[length=2.2mm]},thick}]
\node[box] (a) {Học adapter\\(như V2)};
\node[box,right=0.5cm of a] (b) {STS: dời lớp cũ\\(như V2)};
\node[nb,right=0.5cm of b] (c) {\textbf{1.} Hình dạng\\mượn lớp gần nhất};
\node[nb,right=0.5cm of c] (d) {\textbf{2.} Nới rộng\\$s=2$};
\node[nb,below=0.8cm of d] (e) {\textbf{3.} CA $\times4$,\\lấy trung bình};
\node[box,left=0.5cm of e] (f) {Phân loại\\(một head)};
\draw[arr](a)--(b); \draw[arr](b)--(c); \draw[arr](c)--(d); \draw[arr](d)--(e); \draw[arr](e)--(f);
\end{tikzpicture}
\caption{Một bước của HRC. Ô xanh lá là phần mới so với V2.}
\end{figure}

\subsection{Thành phần 1: mượn hình dạng của lớp gần nhất}
Các lớp giống nhau (ví dụ hai giống chó) có đám đặc trưng hình dạng tương tự. Ta lấy một ``hình dạng tham khảo'' $P_c$
là trung bình có trọng số covariance của các lớp khác, ưu tiên lớp có mean gần $\bmu_c$ nhất \cite{freelunch}:
\begin{equation}
P_c=\sum_{j\ne c}w_{cj}S_j,\qquad w_{cj}=\frac{\exp(\tau\cos(\bmu_c,\bmu_j))}{\sum_{k\ne c}\exp(\tau\cos(\bmu_c,\bmu_k))}.
\end{equation}
Với $\tau=256$, trọng số gần như dồn hết cho lớp gần nhất. Rồi trộn với covariance của chính lớp đó, như thể ta có
thêm $n_0$ ``ảnh giả'' mang hình dạng tham khảo:
\begin{equation}
S_c=\frac{n_0P_c+(n_c-1)C_c}{n_0+n_c-1}+10^{-3}I,\qquad n_0=20 .
\end{equation}
Với 30 ảnh, hình dạng tham khảo chiếm khoảng $40\%$. Các hướng mà 30 ảnh không thấy giờ được lấp bằng hình dạng của
lớp lân cận. Khi $n_0=0$ ta trở lại đúng V2.

\subsection{Thành phần 2: nới rộng đám mây}
Khi sinh dữ liệu giả cho CA, mọi lớp (cũ và mới) đều dùng $\mathcal{N}(\kappa_c\bmu_c,\ s\,S_c)$ với \textbf{cùng}
$s=2$. Hệ số $\kappa_c$ (co mean lớp cũ về gốc) giữ nguyên như RSIAT. Chúng tôi đã thử cho mỗi lớp một $s$ riêng: kết
quả tệ hơn, vì lớp nào có đám rộng hơn sẽ ``chiếm'' vùng của lớp khác. $s$ phải là một hệ số chung.

\subsection{Thành phần 3: lấy trung bình bốn lần CA}
CA huấn luyện trên điểm bốc ngẫu nhiên, nên mỗi lần chạy cho kết quả khác nhau một chút (dao động $\pm0.3$--$0.5$
điểm). HRC chạy CA bốn lần từ cùng điểm xuất phát với bốn lần bốc ngẫu nhiên khác nhau, rồi lấy trung bình hướng của
các vector trọng số. Lần đầu dùng đúng chuỗi số ngẫu nhiên như V2, ba lần sau dùng chuỗi riêng và khôi phục lại, nên
quỹ đạo huấn luyện của các bước sau không đổi.

\subsection{Lớp cũ}
Hình dạng $S_c$ của lớp cũ được STS vận chuyển như V2: $S_c\leftarrow A^\top S_cA$. Không thêm thành phần bất định nào.

% =====================================================================================================================
\section{Thực nghiệm}

\paragraph{Thiết lập.} ImageNet-A: 200 lớp, 10 bước $\times$ 20 lớp, khoảng 30 ảnh huấn luyện mỗi lớp. ViT-B/16
ImageNet-21k với adapter, siêu tham số giống RSIAT gốc. Seed 1993. Đo độ chính xác trên mọi lớp đã gặp sau mỗi bước.
``Avg'' là trung bình 10 bước, ``Last'' là sau bước cuối. Tập kiểm tra nhỏ: một ảnh tương đương $0.06$--$0.3$ điểm.

\begin{table}[H]
\centering\small
\caption{ImageNet-A, seed 1993, một bộ phân loại CA. In đậm: cao nhất mỗi cột.}
\setlength{\tabcolsep}{3.5pt}
\begin{tabular}{l*{10}{c}cc}
\toprule
Bước & 0 & 1 & 2 & 3 & 4 & 5 & 6 & 7 & 8 & 9 & Avg & Last\\
\midrule
RSIAT gốc & 86.29 & 82.50 & 79.62 & \textbf{78.65} & 75.51 & 73.23 & 70.56 & 68.59 & 66.88 & 65.77 & 74.76 & 65.77\\
V2 & 86.29 & 82.22 & 79.62 & 78.51 & 75.62 & \textbf{74.05} & 71.19 & 69.07 & 67.46 & 66.29 & 75.03 & 66.29\\
HRC (1 lần CA) & 86.29 & 83.89 & 80.88 & 78.37 & \textbf{75.86} & 73.64 & \textbf{71.81} & 69.23 & 67.46 & 65.83 & 75.33 & 65.83\\
\textbf{HRC} & 86.29 & \textbf{84.44} & \textbf{81.09} & 78.37 & 75.74 & 73.95 & 71.72 & \textbf{69.55} & \textbf{67.53} & \textbf{66.56} & \textbf{75.52} & \textbf{66.56}\\
\midrule
HRC $-$ V2 & 0.00 & $+2.22$ & $+1.47$ & $-0.14$ & $+0.12$ & $-0.10$ & $+0.53$ & $+0.48$ & $+0.07$ & $+0.27$ & $+0.49$ & $+0.27$\\
HRC $-$ gốc & 0.00 & $+1.94$ & $+1.47$ & $-0.28$ & $+0.23$ & $+0.72$ & $+1.16$ & $+0.96$ & $+0.65$ & $+0.79$ & $+0.76$ & $+0.79$\\
\bottomrule
\end{tabular}
\end{table}

\begin{figure}[H]
\centering
\begin{tikzpicture}
\begin{axis}[width=0.92\linewidth,height=6.2cm,xlabel={Bước},ylabel={Độ chính xác (\%)},xtick={0,...,9},
 ymin=64,ymax=88,grid=major,grid style={black!12},legend style={at={(0.02,0.03)},anchor=south west,font=\small}]
\addplot[cOrig,thick,mark=o] coordinates {(0,86.29)(1,82.50)(2,79.62)(3,78.65)(4,75.51)(5,73.23)(6,70.56)(7,68.59)(8,66.88)(9,65.77)};
\addlegendentry{RSIAT gốc}
\addplot[cV2,thick,mark=square] coordinates {(0,86.29)(1,82.22)(2,79.62)(3,78.51)(4,75.62)(5,74.05)(6,71.19)(7,69.07)(8,67.46)(9,66.29)};
\addlegendentry{V2}
\addplot[cH4,very thick,mark=*] coordinates {(0,86.29)(1,84.44)(2,81.09)(3,78.37)(4,75.74)(5,73.95)(6,71.72)(7,69.55)(8,67.53)(9,66.56)};
\addlegendentry{HRC}
\end{axis}
\end{tikzpicture}
\caption{Bước 0 giống nhau vì chưa có CA. HRC tăng mạnh ở bước 1--2 và giữ lợi thế nhỏ đến bước cuối.}
\end{figure}

\begin{figure}[H]
\centering
\begin{tikzpicture}
\begin{axis}[ybar,bar width=7pt,width=0.92\linewidth,height=5cm,xlabel={Bước},ylabel={HRC $-$ V2 (điểm)},
 xtick={1,...,9},ymin=-0.5,ymax=2.5,grid=major,grid style={black!12}]
\addplot[fill=cH4!70] coordinates {(1,2.22)(2,1.47)(3,-0.14)(4,0.12)(5,-0.10)(6,0.53)(7,0.48)(8,0.07)(9,0.27)};
\end{axis}
\end{tikzpicture}
\caption{Chênh lệch theo bước giữa HRC và V2.}
\end{figure}

\paragraph{Đóng góp của từng phần.} Mượn hình dạng và nới rộng ($s=2$) cho $+0.30$ Avg nhưng Last còn $-0.46$ (một lần
CA). Lấy trung bình bốn lần CA kéo Last lên $66.56$ và Avg lên $75.52$. Mô phỏng ngoại tuyến (cùng backbone, lặp lại
nhiều lần bốc ngẫu nhiên) cho thấy riêng phần mượn hình dạng và nới rộng mang lại khoảng $+0.3$ đến $+0.45$ Avg; lấy
trung bình chủ yếu giảm dao động và cải thiện Last.

\paragraph{Seed khác.} Một lần chạy seed 1994 (dừng ở bước 5) cho HRC một CA hơn V2 ở mọi bước:
$+1.88,+0.67,+1.30,+1.29,+1.00$ ở bước 1--5.

% =====================================================================================================================
\section{Những gì đã thử nhưng không thành công}
\begin{itemize}
\item \textbf{Chọn tham số ``trung thực''.} Ý định ban đầu của HRC: huấn luyện một adapter phụ trên một nửa ảnh mỗi
      lớp, dùng nửa còn lại (mô hình chưa thấy) để chọn $(n_0,\tau,s)$. Mọi tiêu chí (log-loss, log-loss sau hiệu chỉnh
      nhiệt độ, tỉ lệ lỗi, log-likelihood Gauss) đều chọn ra biên của lưới hoặc chọn theo nhiễu. Lý do gốc: nửa giữ lại
      chỉ có lớp mới, nên không quan sát được việc lớp mới bị nhầm thành lớp cũ --- đúng loại lỗi chiếm ưu thế.
\item \textbf{Bất định tích luỹ (V3)}: bỏ hay giữ đều không giúp CA.
\item \textbf{$s$ riêng cho từng lớp} hoặc tách $s$ lớp cũ/mới, bỏ $\kappa$: không tốt hơn.
\item \textbf{Cộng ``khoảng chênh tổng quát hoá''} (phần covariance ảnh thật rộng hơn ảnh huấn luyện): khi bỏ rò rỉ
      tập kiểm tra, lợi ích gần bằng 0.
\end{itemize}

% =====================================================================================================================
\section{Giới hạn cần nói rõ}
\begin{itemize}
\item Bộ tham số $(n_0=20,\tau=256,s=2)$ được chọn bằng một phép quét ngoại tuyến \textbf{có dùng tập kiểm tra} của
      seed 1993. Vì vậy con số trên seed 1993 là lạc quan; cần kiểm chứng trên seed và dữ liệu khác (seed 1994 dang dở
      là dấu hiệu tích cực đầu tiên).
\item Mức tăng nhỏ ($+0.49/+0.27$ so với V2), cỡ vài ảnh trên tập kiểm tra.
\item HRC một bộ phân loại vẫn thấp hơn V2 có tổ hợp ba bộ phân loại ($75.63/66.69$).
\item Trần lý thuyết của việc chỉ sửa ``đám mây'' (với backbone cố định) đo được khoảng $+0.4$--$0.5$ Avg; muốn tăng
      nhiều hơn phải cải thiện chính bộ trích đặc trưng.
\end{itemize}

\section{Kết luận}
Trong học tăng dần không lưu ảnh, bộ phân loại chỉ tốt bằng các ``đám mây'' dùng để nhớ lại lớp cũ. HRC sửa hình dạng
(mượn từ lớp gần nhất), độ rộng (nới gấp đôi cho mọi lớp) và giảm may rủi của CA (trung bình bốn lần). Chỉ với một bộ
phân loại, trên ImageNet-A seed 1993, HRC đạt $75.52/66.56$, cao hơn V2 ($75.03/66.29$) và RSIAT gốc
($74.76/65.77$). Mã: \texttt{models/hrc.py}, cấu hình \texttt{exps/hrc/imageneta\_hrcfix\_avg4\_s1993.json}.

\begin{thebibliography}{99}
\bibitem{ssiat} Y. Tan et al. Semantically-Shifted Incremental Adapter-Tuning is A Continual ViTransformer. \emph{CVPR}, 2024.
\bibitem{slca} G. Zhang et al. SLCA: Slow Learner with Classifier Alignment for Continual Learning on a Pre-trained Model. \emph{ICCV}, 2023.
\bibitem{sdc} L. Yu et al. Semantic Drift Compensation for Class-Incremental Learning. \emph{CVPR}, 2020.
\bibitem{cosface} H. Wang et al. CosFace: Large Margin Cosine Loss for Deep Face Recognition. \emph{CVPR}, 2018.
\bibitem{freelunch} D. Goswami et al. FeCAM / Calibrating Higher-Order Statistics for Few-Shot Class-Incremental Learning with Pre-trained Vision Transformers. 2023--2024.
\bibitem{devroye} L. Devroye, L. Györfi, G. Lugosi. \emph{A Probabilistic Theory of Pattern Recognition}. Springer, 1996.
\bibitem{nc} V. Papyan, X. Y. Han, D. Donoho. Prevalence of Neural Collapse during the Terminal Phase of Deep Learning Training. \emph{PNAS}, 2020.
\bibitem{efron} B. Efron. The Estimation of Prediction Error: Covariance Penalties and Cross-Validation. \emph{JASA}, 2004.
\end{thebibliography}
\end{document}
